\chapter{Experiment 2: Corrected Nested Tensor Train and Tucker Benchmark}
\label{ch:v2}

\section{Corrections and representation changes}

V2 is the first fully corrected tensor benchmark. It verifies the 19-channel
EEGLAB row order for all 69 controlled-cohort participants, applies the 0.5 s
boundary guard, and retains 11,942 windows. It also replaces relative,
non-negative power with signed log absolute PSD so that centered increases and
decreases remain available to an orthogonal low-rank model.

Each retained window is represented as
$\mathcal{Z}_w\in\mathbb{R}^{19\times59\times3}$: channel $\times$ frequency
$\times$ local segment. A separate subject tensor stores mean, median, and
standard deviation logPSD, also of shape $19\times59\times3$. The local segment
axis is not described as shared event time.

\section{Spectral controls}

The V2 control vector has 418 features. For delta (1--4 Hz), theta (4--8 Hz),
alpha (8--13 Hz), and beta (13--30 Hz), it records per-channel mean and standard
deviation of log absolute band power and relative band power. It also records
per-channel mean and standard deviation of spectral entropy, 7--13 Hz alpha
peak, and log slow-to-fast power ratio. Age, sex, MMSE, duration, window count,
and diagnosis-derived statistics are absent.

\section{Inductive TT and Tucker projections}

The \code{TTProjection.fit()} implementation centers its training tensor,
moves observations to the last mode, and computes leading left vectors from
$AA^T$ for each unfolding. Signs are fixed by the largest-magnitude element for
reproducible cores. The contracted cores form a common orthonormal basis; the
code records orthogonality and centered reconstruction error.

Window-level bases use 24 chronologically spread calibration windows per
training participant. For every participant, all projected windows are pooled
using four statistics per latent component: mean, standard deviation, median,
and interquartile range. The decomposition basis is trained only from the
relevant training participants; calibration provenance is saved.

Tucker uses training-only multilinear PCA/HOSVD ranks $(6,8,3)$. Graph TT first
applies a fixed coordinate-derived smoother. Each electrode connects to its four
nearest coordinate neighbors with Gaussian distance weights. For normalized
Laplacian $L=I-D^{-1/2}AD^{-1/2}$ and strength $\eta=0.5$, the smoother is

\begin{equation}
S=V\operatorname{diag}\left((1+\eta\lambda_i)^{-1/2}\right)V^T,
\end{equation}

where $L=V\operatorname{diag}(\lambda_i)V^T$. This is deterministic anatomical
pre-smoothing, not the preliminary graph penalty in Eq.~\ref{eq:graphcp}.

\section{Candidate families and models}

Six representation families are compared:

\begin{enumerate}
  \item subject-level TT at ranks 8, 16, and 32;
  \item pooled window TT at ranks 8, 16, and 32;
  \item graph-smoothed pooled window TT at ranks 8, 16, and 32;
  \item pooled window Tucker with training-only selection of 16, 32, or 64
  latent features;
  \item TT plus 418 spectral controls with TT ranks 8, 16, or 32 and selection
  of 32 combined features; and
  \item spectral controls with selection of 16, 32, or 64 features.
\end{enumerate}

Classifier candidates are logistic regression with $C\in\{0.1,1,10\}$, linear
SVM with $C\in\{0.1,1\}$, RBF SVM with $C\in\{1,10\}$, and a fixed shallow
histogram gradient booster. Feature selection uses ANOVA $F$ scores; every
selector and scaler is inside the fitted pipeline.

\section{Nested subject cross-validation}

The outer evaluation is three repetitions of stratified five-fold participant
CV (seed 5040). Each outer training set contains an independent three-fold inner
CV (seed 8000 plus outer index). Inner out-of-fold balanced accuracy, then macro
F1, selects rank, representation, feature count, and classifier. Exact ties use
a deterministic lower-complexity ordering. This design addresses the optimistic
bias produced when the same cross-validation result both selects and evaluates a
model \citep{varma2006nested}.

\begin{figure}[!htbp]
\centering
\begin{tikzpicture}[node distance=9mm and 9mm,
block/.append style={text width=63mm,minimum height=17mm},
process/.append style={text width=63mm,minimum height=17mm}]
\node[block,text width=115mm] (people) at (0,0) {69 participants in V2; 88 in V3 / V4\\Outer CV: 3 repetitions of 5 folds};
\node[process] (train) at (-3.8,-2.4) {Outer training people\\Inner 3-fold CV chooses the rank,\\features and classifier configuration};
\node[block] (test) at (3.8,-2.4) {Outer test people\\All windows from these people\\are withheld from fitting and tuning};
\node[process,below=of train] (refit) {Refit the selected complete pipeline\\Basis, selector, scaler, weights and\\classifier use outer training people};
\node[block,below=of test] (predict) {Predict outer test people\\One prediction per person per repeat\\Pool complete-repeat subject metrics};
\draw[arrow] (people.south) -- ++(0,-0.3) -| (train.north);
\draw[arrow] (people.south) -- ++(0,-0.3) -| (test.north);
\draw[arrow] (train)--(refit);
\draw[arrow] (test)--(predict);
\draw[arrow] (refit)--(predict);
\end{tikzpicture}
\caption{Nested participant-level training, validation and test preparation in V2--V4. Inner validation selects a configuration; outer test labels enter only the final scoring step.}
\label{fig:nested-cv}
\end{figure}


The \emph{primary} V2 endpoint is the outer score of the entire inner-selected
pipeline. Per-family outer results are secondary method comparisons. SVM class
labels come from \code{predict()}, while its decision scores are stored as
uncalibrated decision functions rather than probabilities.

\section{Results and interpretation}

\begin{table}[htbp]
\centering
\begin{tabular}{lrr}
\toprule
\textbf{V2 family} & \textbf{Mean BA} & \textbf{Repeat SD}\\
\midrule
Inner-selected complete pipeline (primary) & 47.3\% & --\\
Window Tucker / HOSVD & 56.0\% & 3.3 pp\\
TT plus spectral controls & 54.1\% & --\\
Graph-smoothed window TT & 53.1\% & --\\
Spectral controls & 51.7\% & --\\
Window TT & 51.2\% & --\\
Subject TT & 50.2\% & 3.0 pp\\
\bottomrule
\end{tabular}
\caption{Corrected V2 nested subject-level results. Dashes indicate values not emphasized in the consolidated record.}
\label{tab:v2-results}
\end{table}

The strongest observed family is Tucker/HOSVD at 56.0\% mean balanced accuracy,
with mean recalls of 59.4\% \AD{}, 50.7\% \FTD{}, and 58.0\% \CN{}. However,
the complete automatic-selection procedure reaches only 47.3\%. The gap shows
selection variance in a small cohort: selecting among many near-tied pipelines
inside roughly 55 training participants can fail to identify the family that
looks strongest after aggregation.

\section{Verification}

Numerical tests compare TT reconstruction against TensorLy, verify orthogonality,
recover full-rank synthetic tensors, and check batch-invariant transform output.
Post-run verification replays 105 saved classifier pipelines, 60 decomposition
sets, 45 inner splits, and 15 outer splits. Held-out features are recomputed from
saved training bases and must reproduce predictions and scores. Input hashes and
source transforms are checked; no test participant enters basis calibration.

